\documentclass[lang=en,color=blue,mode=fancy,titlestyle=hang]{elegantbook}

% ---- Title Information ----
% with lang=en the class does not load ctex, and any Chinese in the body renders as blank.
% If you really need Chinese text, add \usepackage{xeCJK} and \setCJKmainfont{...}.
\title{Title of the Tutorial}
\author{Oct0pu5}
\date{Jan 1st, 2027}
\version{1.0}
\extrainfo{``A short quotation, printed at the foot of the title page.''}
\cover{cover.pdf}

% ---- Additional Packages ----
% Only what you actually use. The class already loads amsmath, graphicx, hyperref, xcolor,
% listings and more, so do not repeat those.
\usepackage{array}
\usepackage{booktabs}
\usepackage{float}
\usepackage{wrapfig}
\usepackage{tabularx}
\usepackage{enumitem}
\usepackage{tikz}

% ---- Math Operators ----
\DeclareMathOperator*{\argmin}{arg\,min}
\DeclareMathOperator*{\argmax}{arg\,max}

% ---- Figure captions ----
% Suppresses the "Figure N:" prefix. No figure in these tutorials is ever cross-referenced,
% so the number carries no information. This line is easy to lose when a preamble is
% rewritten, because it is not attached to any \usepackage.
\usepackage{caption}
\captionsetup[figure]{labelformat=empty}

% ---- Table of contents ----
% Entries at body size, opened up with leading. Enlarging the font instead makes the page
% look bulky; spacing fills it without shouting.
\setlength{\cftbeforechapskip}{1.1em}
\setlength{\cftbeforesecskip}{0.5em}
\setlength{\cftbeforesubsecskip}{0.15em}

\begin{document}

\maketitle

\frontmatter
\tableofcontents

\mainmatter

% ============================================================
\chapter*{Preface}
\addcontentsline{toc}{chapter}{Preface}

% Wider leading for the Preface. Note that the class sets \baselinestretch to 1.2, not 1,
% so the reset below has to be 1.2. Writing 1 tightens the whole book.
\linespread{1.55}\selectfont

Write why you took this subject on, whose course or book the tutorial follows, who it is
for, and one piece of advice for the reader.

\linespread{1.2}\selectfont

% ============================================================
\chapter{First Chapter}

\subsection*{Overview}

% Three short paragraphs. Measure before you write, rather than after: the target is this
% document's own average sentence length, not another tutorial's. To measure it:
%
%   detex() { sed -e 's/\\[a-zA-Z]*//g' -e 's/[${}\\]//g' | tr '\n' ' '; }
%   w=$(detex < file.tex | wc -w); s=$(detex < file.tex | grep -oE '[.!?] ' | wc -l)
%   echo $((w / s))
%
% Around eighteen words per sentence is the house style. Em-dashes are rare, semicolons
% rarer still.
%
%   Paragraph 1: what problem the chapter solves, and why it sits here.
%   Paragraph 2: what it actually contains.
%   Paragraph 3: how it connects to the chapters around it.

Say what problem this chapter solves and why it comes here.

Say what the chapter actually contains.

Say how it connects to the chapters around it.

\section{First Section}

Body text.

\subsection{A Subsection}

Body text.

% ============================================================
\chapter*{Afterwords}
\addcontentsline{toc}{chapter}{Afterwords}

\linespread{1.55}\selectfont

Write what writing the tutorial taught you, what you could not cover, and any advice for a
reader who wants to write their own notes.

\begin{flushright}
  Oct0pu5\\
  Jan 1st, 2027
\end{flushright}

\end{document}
